\documentclass[tikz,border=10pt]{standalone}
\usepackage{tikz}
\usepackage{amsmath}
\usepackage{amssymb}
\usepackage{pgfplots}
\pgfplotsset{compat=1.18}
\usepackage{ctex}
\usetikzlibrary{calc, positioning, arrows.meta, decorations.pathreplacing, angles, quotes}

\begin{document}
\begin{tikzpicture}[scale=1.1]

  % 背景网格（淡）
  \draw[gray!20, thin] (-0.5,-1.2) grid (5.5,4.2);
  \draw[->,gray!50, thin] (-0.5,0) -- (5.8,0) node[right,font=\scriptsize] {$x$};
  \draw[->,gray!50, thin] (0,-1.2) -- (0,4.5) node[above,font=\scriptsize] {$y$};

  % 定义坐标
  % u = (2, 4), v = (5, 1)
  % proj of u onto v: scalar = (u·v)/(v·v) = (10+4)/(25+1) = 14/26 ≈ 0.5385
  % proj vector F = 0.5385*(5,1) = (2.692, 0.538)
  \coordinate (O) at (0,0);
  \coordinate (U) at (2,4);
  \coordinate (V) at (5,1);
  \coordinate (F) at (2.692, 0.538);

  % 向量 v（底色，最先画）
  \draw[->,thick,teal!80!black,>=Stealth] (O) -- (V)
    node[below right,font=\small] {$\mathbf{v}$};

  % 向量 u
  \draw[->,thick,blue!80!black,>=Stealth] (O) -- (U)
    node[above left,font=\small] {$\mathbf{u}$};

  % 投影虚线（从 U 到 F）
  \draw[densely dashed, gray!60, thick] (U) -- (F);

  % 投影向量（红色实箭头）
  \draw[->,thick,red!80!black,>=Stealth] (O) -- (F)
    node[below,font=\small,yshift=-6pt]
      {$\mathrm{proj}_{\mathbf{v}}\mathbf{u}$};

  % 垂直分量 u - proj（绿色）
  \draw[->,thick,green!60!black,>=Stealth] (F) -- (U)
    node[right,font=\small, xshift=4pt]
      {$\mathbf{u}-\mathrm{proj}_{\mathbf{v}}\mathbf{u}$};

  % 直角标记
  \draw[gray!70, thin]
    ($(F)!0.15!(U)$) --
    ($(F)!0.15!(U) + ($(F)!0.15!(V)$) - (F)$) --
    ($(F)!0.15!(V)$);

  % 点标注
  \fill (O) circle (1.5pt) node[below left,font=\scriptsize] {$O$};
  \fill (F) circle (1.5pt);

  % 公式框（白底灰边，放在图外下方）
  \node[font=\small, align=center, fill=white, draw=gray!50, thin,
        rounded corners=3pt, inner sep=7pt]
    at (2.8, -1.8) {
      $\displaystyle\mathrm{proj}_{\mathbf{v}}\mathbf{u}
        = \frac{\mathbf{u}\cdot\mathbf{v}}{\mathbf{v}\cdot\mathbf{v}}\,\mathbf{v}$
    };

  % 标题
  \node[font=\large\bfseries] at (2.5, 5.2) {内积与投影};

\end{tikzpicture}
\end{document}
